% !TeX root = ../../Book3.tex
\section{无穷乘积}
\label{b3:sec:2001}

有限个零点可以用多项式乘积实现。无穷个零点要求控制乘积尾部，且必须确认极限除了指定因子之外不会产生额外零点。

\subsection{收敛判据}
\label{b3:subsec:200101}

若 \(\sum_n|u_n|<\infty\)，则 \(\prod_n(1+u_n)\) 收敛；当所有因子非零时极限也非零。证明取尾部 \(|u_n|\le1/2\)，使用 \(\log(1+u_n)=\sum_{k\ge1}(-1)^{k+1}u_n^k/k\) 的固定小圆分支，得到 \(|\log(1+u_n)|\le2|u_n|\)。尾部乘积是收敛对数和的指数，故非零；有限头部另行相乘。

不控制对数和，仅有 \(u_n\to0\) 不足以保证收敛：\(\prod(1+1/n)\) 发散。也不能把“收敛”一律解释为非零；\(\prod(1-1/n)\) 从第二项起趋零。

\subsection{绝对收敛}
\label{b3:subsec:200102}

绝对收敛的乘积允许重排，其原因是小因子的对数绝对可和。无穷乘积的“绝对收敛”通常指 \(\sum|u_n|<\infty\)，不是对乘积模再作一次级数。

若有某个因子为零，整乘积为零，此时对数表述只用于去掉该因子后的尾部。后面在函数论中，零点附近总会把产生零点的有限因子与无零点的尾部拆开。

\subsection{正规收敛}
\label{b3:subsec:200103}

\begin{proposition}\label{b3:prop:20-normal-product}
若全纯函数 \(u_n\) 在区域 \(D\) 上满足对每个紧集 \(K\Subset D\)，\(\sum_n\sup_K|u_n|<\infty\)，则 \(\prod(1+u_n)\) 局部一致收敛到全纯函数；在无因子零点的地方，极限非零。
\end{proposition}
\begin{proof}
固定紧集附近稍大的紧集，尾部因子一致落在 \(|u|\le1/2\)，对数级数由上述估计正规收敛，其和全纯，指数给全纯非零尾部。乘上有限头部，得到局部一致极限。不同紧集的尾部表示来自同一有限乘积序列，故彼此一致。
\end{proof}
局部一致性保证可逐项微分对数：无零点的紧子集上 \(P'/P=\sum u_n'/(1+u_n)\)，导数收敛由全纯局部一致收敛定理给出。

\subsection{零点与乘积}
\label{b3:subsec:200104}

在固定点附近，若只有有限个因子为零，就把它们提出；余下尾部由正规对数和给非零函数。因此极限的零点阶数恰为有限零因子的阶数和，没有额外零点。

如果零因子在内部无限聚集，极限可能恒零，而指定离散零点的任务本身就与恒等定理冲突。本章的零点序列始终要求无有限聚点。


对乘积尾部还可给出误差界。若 \(S_N(K)=\sum_{n>N}\sup_K|u_n|\)，则由有限乘积展开
\[
\left|\prod_{n=N+1}^{M}(1+u_n)-1\right|\le
\prod_{n=N+1}^{M}(1+|u_n|)-1\le e^{S_N(K)}-1.
\]
令 \(M\to\infty\)，这个界仍成立。它直接控制尾部相对误差，而不仅是存在一个极限。有限头部在紧集的最大模有界，所以全乘积的绝对误差亦受同一尾和控制。若需要数值计算，先选 \(N\) 使尾和小，再单独计算头部；没有尾部界的有限截断不能证明函数恒等式。

\ProseParagraphBreak

以 \(u_n(z)=z/n^2\) 为例，\(|z|\le R\) 时 \(S_N\le R/N\)。只要没有碰到有限头部的零点，截断相对误差不超过 \(e^{R/N}-1\)。若零点是当前计算目标，必须提出相应因子后再谈相对误差，因为在零值处除以极限没有意义。相同估计还说明局部一致性不能改成全平面一致性：随着 \(R\) 增大，所需 \(N\) 也增大。

\ProseParagraphBreak

条件收敛的乘积可能依赖顺序。例如实数列 \(u_n=(-1)^n/n\) 从 \(n=2\) 起满足 \(\sum u_n\) 条件收敛、\(\sum u_n^2<\infty\)。由
\(\log(1+u_n)=u_n+O(u_n^2)\)，原顺序乘积收敛为正数，但正负项重排能改变一次和，从而改变乘积。这里所有因子甚至都正；失败来自绝对收敛不足，而不是复杂的复数分支。典范因子就是有目的地消除这些可能依赖顺序的低阶和。


\ProseParagraphBreak

正规收敛还允许检验导数，但须区分乘积与对数的适用位置。若紧集避开所有零点，取稍大的无零点邻域，并在那里对尾部统一选择对数，则 Cauchy 公式给对数级数各阶导数的局部一致收敛。若紧集包含一个零点，先分离该点的有限零因子，然后只对剩下的非零部分取对数。若直接写 \(P'/P=\sum u_n'/(1+u_n)\)，左右都是亚纯函数；在零点处谈全纯一致收敛并无意义。

\ProseParagraphBreak

例如 \(P(z)=\prod_{n\ge1}(1+z/n^2)\) 在原点附近的对数为
\[
\log P(z)=\sum_{k\ge1}\frac{(-1)^{k+1}z^k}{k}\sum_{n\ge1}n^{-2k},\qquad |z|<1.
\]
先取 \(|z|\le r<1\)，双重级数的绝对值由 \(\sum_n\sum_k r^k/n^{2k}\) 控制，后者不超过 \(r(1-r)^{-1}\sum_n n^{-2}\)。于是取原点一次导数得到 \(P'(0)=\sum n^{-2}\)，二次导数则为 \((\sum n^{-2})^2-\sum n^{-4}\)。对数系数记录单因子的贡献，函数系数还包含因子之间的交叉项；把这两类系数混在一起会漏掉平方项。

\ProseParagraphBreak

无穷乘积定理不能仅由各因子趋于一替代。常数因子 \(1+1/n\) 虽趋于一，其前 \(N\) 项的乘积为 \(N+1\)，仍发散。常数因子 \(1-1/n\) 从第二项开始乘积为 \(1/N\)，极限恒为零，不能用来断言“除有限个零因子以外尾部非零”。这两个计算说明，可和的对数尾项同时承担有限性与非消失性两项责任。判别时应检查它控制哪一项，不能只检查逐因子的极限。
